\documentclass[10pt,conference]{IEEEtran}

% ── Packages ──────────────────────────────────────────────────────
\usepackage[utf8]{inputenc}
\usepackage{times}
\usepackage{cite}
\usepackage{amsmath,amssymb,amsfonts}
\usepackage{graphicx}
\usepackage{textcomp}
\usepackage{xcolor}
\usepackage{url}
\usepackage{listings}
\usepackage{booktabs}
\usepackage{array}
\usepackage{tabularx}
\usepackage{multirow}
\usepackage{hyperref}

% ── Code listing style ─────────────────────────────────────────────
\lstset{
  basicstyle=\ttfamily\footnotesize,
  breaklines=true,
  frame=single,
  xleftmargin=4pt,
  xrightmargin=4pt,
  captionpos=b
}

% ── Document ───────────────────────────────────────────────────────
\begin{document}

\title{SecureShift: A DevSecOps Pipeline with\\
Automated Multi-Layer Security Scanning\\
and Real-Time Threat Intelligence Dashboard}

\author{
  \IEEEauthorblockN{Neel Patel}
  \IEEEauthorblockA{
    Dept. of Computer Engineering\\
    San José State University\\
    San José, CA, USA\\
    neel.patel@sjsu.edu
  }
  \and
  \IEEEauthorblockN{Author Two}
  \IEEEauthorblockA{
    Dept. of Computer Engineering\\
    San José State University\\
    San José, CA, USA\\
    author2@sjsu.edu
  }
  \and
  \IEEEauthorblockN{Author Three}
  \IEEEauthorblockA{
    Dept. of Computer Engineering\\
    San José State University\\
    San José, CA, USA\\
    author3@sjsu.edu
  }
}

\maketitle

% ── Abstract ───────────────────────────────────────────────────────
\begin{abstract}
Modern software delivery increasingly demands security assurance at every stage
of the development lifecycle rather than as a post-deployment afterthought.
This paper presents \emph{SecureShift}, a fully automated DevSecOps pipeline
that integrates four complementary security scanning disciplines—dependency
auditing, static application security testing (SAST), container image
scanning, and dynamic application security testing (DAST)—into a single
GitHub Actions CI/CD workflow.
To demonstrate measurable security improvement, the pipeline is evaluated
against a purpose-built, intentionally vulnerable Node.js/Express REST API
containing thirteen documented weaknesses spanning OWASP Top~10 categories
including SQL injection, stored cross-site scripting, command injection, path
traversal, insecure direct object reference, and hard-coded credentials.
Scan outputs from npm audit, SonarCloud, Aqua Trivy, and OWASP ZAP are
normalized and published to a React-based Security Intelligence Dashboard
hosted on Render, giving development teams real-time visibility into
their security posture.
A remediated version of the same API (app-fixed) then demonstrates that the
pipeline can correctly confirm a clean bill of health once all vulnerabilities
are addressed.
Our evaluation shows that the four-tool combination surfaces all thirteen
intentional vulnerabilities with zero configuration beyond a single YAML
workflow file, confirming the feasibility of shift-left security for teams
with limited security expertise.
\end{abstract}

\begin{IEEEkeywords}
DevSecOps, CI/CD, SAST, DAST, container security, OWASP, GitHub Actions,
security dashboard, vulnerability scanning, shift-left security
\end{IEEEkeywords}

% ── I. INTRODUCTION ────────────────────────────────────────────────
\section{Introduction}

The traditional software development model treats security as a final gate
before production, commonly referred to as ``bolt-on'' security.
This approach produces expensive rework: the National Institute of Standards
and Technology (NIST) estimates that defects discovered in production can cost
up to thirty times more to remediate than those caught during
development~\cite{nist2002cost}.
The DevSecOps movement addresses this systemic problem by embedding security
controls directly into Continuous Integration/Continuous Deployment (CI/CD)
pipelines so that every code commit is automatically evaluated against a
portfolio of security checks~\cite{devsecops2019}.

Despite growing industry adoption, many development teams still face
significant barriers to entry.
Security tooling is fragmented across vendors, output formats are inconsistent,
and interpreting scanner results requires domain expertise that small
engineering teams rarely possess.
This work addresses those barriers with a practical, end-to-end reference
implementation that any team can deploy with a single GitHub repository.

\subsection{Contributions}

The primary contributions of this paper are:

\begin{enumerate}
  \item A fully automated, six-stage GitHub Actions DevSecOps pipeline
        combining dependency scanning (npm audit), SAST (SonarCloud), container
        scanning (Trivy), and DAST (OWASP ZAP) in a single declarative YAML
        workflow.

  \item A purpose-built intentionally vulnerable Node.js/Express API
        (``VulnAPI'') containing thirteen documented vulnerabilities across seven
        OWASP Top~10 categories, designed as a reproducible evaluation target.

  \item A React-based Security Intelligence Dashboard that aggregates,
        normalizes, and visualizes findings from all four scanners in real time
        and is automatically redeployed to Render on every pipeline run.

  \item A remediated version of VulnAPI (``app-fixed'') demonstrating the
        end-to-end workflow from vulnerability detection to resolution.
\end{enumerate}

\subsection{Paper Organization}

The remainder of this paper is organized as follows.
Section~II reviews related work.
Section~III describes the overall system architecture.
Section~IV details the vulnerable application design.
Section~V describes each pipeline stage.
Section~VI presents the security dashboard.
Section~VII covers the remediation strategy.
Section~VIII presents experimental results.
Section~IX discusses limitations and future work.
Section~X concludes the paper.

% ── II. RELATED WORK ───────────────────────────────────────────────
\section{Related Work}

\subsection{DevSecOps and Shift-Left Security}

The term DevSecOps emerged as an extension of the DevOps movement, formally
introduced to emphasize that security should be a shared responsibility across
development, security, and operations teams~\cite{devsecops2019}.
Myrbakken and Colomo-Palacios~\cite{myrbakken2017} conducted an early
practitioner survey finding that CI/CD integration was the single highest-value
control for reducing mean time to remediation.
More recent work by Rajapakse et al.~\cite{rajapakse2022} systematically
reviewed DevSecOps challenges, identifying tool fragmentation and the absence
of unified dashboards as the leading inhibitors of adoption.
Our work directly addresses these gaps.

\subsection{Automated Vulnerability Scanning}

Npm audit, introduced by the Node Package Manager team in 2018, applies
advisory data from the National Vulnerability Database (NVD) to compute
a dependency graph and flag transitive vulnerabilities~\cite{npmaudit}.
SonarCloud extends the SonarQube SAST engine to the cloud, enforcing quality
and security rules across dozens of languages without self-hosting
overhead~\cite{sonarcloud}.
Aqua Trivy performs OS-level and library-level CVE scanning against container
images, providing SARIF output compatible with GitHub's Security
tab~\cite{trivy2021}.
OWASP ZAP remains the most widely adopted open-source DAST tool, capable of
active scanning with authenticated sessions and API schema
support~\cite{zapproxy}.

\subsection{Security Dashboards}

Previous work on security visibility has largely focused on Security Information
and Event Management (SIEM) platforms such as Splunk and Elastic SIEM, which
require significant infrastructure investment~\cite{siem2020}.
Lighter-weight, developer-facing dashboards that surface CI/CD scan data
directly to pull-request authors are less studied.
Perez-Castillo et al.~\cite{perez2021} proposed a developer dashboard embedded
in the IDE; our approach instead targets web-based, always-on visibility
accessible to all stakeholders without IDE access.

% ── III. SYSTEM ARCHITECTURE ───────────────────────────────────────
\section{System Architecture}

\begin{figure}[b]
  \centering
  \setlength{\fboxsep}{6pt}
  \fbox{%
    \begin{minipage}{0.85\columnwidth}
      \centering
      \ttfamily\scriptsize
      \textbf{Developer Push}\\[4pt]
      $\downarrow$\\[2pt]
      \textbf{GitHub Actions Workflow}\\[4pt]
      \begin{tabular}{ccc}
        npm audit & SonarCloud & Docker Build \\
        $\downarrow$ & $\downarrow$ & $\downarrow$ \\
        -- & -- & Trivy + ZAP \\
      \end{tabular}\\[4pt]
      $\downarrow$\\[2pt]
      \textbf{Dashboard Build \& Deploy}\\[4pt]
      $\downarrow$\\[2pt]
      \textbf{Render (Security Dashboard)}
    \end{minipage}%
  }
  \caption{SecureShift system architecture. A code push triggers the GitHub
           Actions pipeline, which runs four scanning jobs, collects
           JSON reports, and deploys the aggregated dashboard to Render.}
  \label{fig:architecture}
\end{figure}

Figure~\ref{fig:architecture} illustrates the high-level architecture.
Three logical tiers interact in the system:

\textbf{Source Tier.} A GitHub repository hosts both the application under
test and the pipeline definition (\texttt{.github/workflows/devsecops.yml}).
Every push to any branch and every pull request targeting \texttt{main} triggers the
pipeline.

\textbf{Pipeline Tier.} Six sequential and parallel GitHub Actions jobs execute
inside ephemeral Ubuntu runners:
(1)~Dependency Scan, (2)~SAST, (3)~Docker Build, (4)~Container Scan,
(5)~DAST, and (6)~Dashboard Build~\&~Deploy.
Each scan job uploads a structured JSON report as a GitHub Actions artifact.
The final job downloads all artifacts, injects pipeline metadata, builds the
React dashboard, and deploys it.

\textbf{Presentation Tier.} The React Security Intelligence Dashboard, hosted
on Render's free-tier static site service at a public HTTPS URL,
fetches the JSON reports at runtime and renders interactive charts and threat
tables.
This design requires no server-side infrastructure beyond Render's free-tier
static hosting.

The system is fully declarative: the entire pipeline is captured in a
single 363-line YAML file, requiring only one repository secret
(\texttt{SONAR\_TOKEN}) and one Render Static Site to be configured.

% ── IV. VULNERABLE APPLICATION (VulnAPI) ───────────────────────────
\section{Vulnerable Application Design}

\subsection{Overview}

VulnAPI is a Node.js~18/Express REST API backed by an in-process SQLite
database.
It exposes five route groups—\texttt{/api/auth}, \texttt{/api/users},
\texttt{/api/products}, \texttt{/api/files}, and \texttt{/api/admin}—and is
containerized with a deliberately insecure Dockerfile.
Table~\ref{tab:vulns} enumerates all thirteen intentional vulnerabilities.

\begin{table}[t]
  \caption{Intentional Vulnerabilities in VulnAPI}
  \label{tab:vulns}
  \centering
  \scriptsize
  \begin{tabular}{@{}llll@{}}
    \toprule
    ID & CWE & Category & Location \\
    \midrule
    VULN-01 & CWE-798 & Hard-coded JWT secret        & \texttt{config.js:12} \\
    VULN-02 & CWE-798 & Hard-coded credentials/keys  & \texttt{config.js:15--21} \\
    VULN-03 & CWE-942 & Wildcard CORS + credentials  & \texttt{server.js:31} \\
    VULN-04 & CWE-89  & SQL injection (login)        & \texttt{routes/auth.js:23} \\
    VULN-05 & CWE-209 & SQL error leaks query text   & \texttt{routes/auth.js:27} \\
    VULN-06 & CWE-639 & IDOR (no ownership check)   & \texttt{routes/users.js:18} \\
    VULN-07 & CWE-312 & Passwords exposed in GET     & \texttt{routes/users.js:9} \\
    VULN-08 & CWE-89  & SQL injection (search)       & \texttt{routes/products.js:14} \\
    VULN-09 & CWE-79  & Stored XSS (unescaped HTML)  & \texttt{routes/products.js:32} \\
    VULN-10 & CWE-22  & Path traversal (file read)  & \texttt{routes/files.js:14} \\
    VULN-11 & CWE-77  & Command injection (ping)    & \texttt{routes/admin.js:18} \\
    VULN-12 & CWE-306 & Unauthenticated admin routes & \texttt{routes/admin.js:6} \\
    VULN-13 & CWE-250 & Root user in Dockerfile      & \texttt{Dockerfile} \\
    \bottomrule
  \end{tabular}
\end{table}

\subsection{Selected Vulnerability Details}

\textbf{VULN-01/02 — Hard-coded Secrets (CWE-798).}
The configuration module (\texttt{config.js}) exports a static JWT signing
secret (\texttt{'supersecret123'}), admin API key, database password, and
a pair of mock AWS access keys.
Hard-coded secrets are trivially discoverable via source code search and
will be flagged by both SonarCloud's \texttt{javascript:S2068} rule and
GitHub's push protection secret scanning.

\textbf{VULN-04 — SQL Injection (CWE-89).}
The login handler constructs queries via string concatenation:
\begin{lstlisting}[language=SQL,caption=Vulnerable login query]
SELECT * FROM users
WHERE username = '${username}'
  AND password = '${password}'
\end{lstlisting}
Supplying \texttt{admin'--} as the username bypasses password verification
entirely.
The same pattern appears in the product search endpoint (VULN-08), enabling
data exfiltration via \texttt{UNION}-based injection.

\textbf{VULN-09 — Stored XSS (CWE-79).}
The product creation endpoint accepts a \texttt{description} field and
stores it verbatim.
When the product list endpoint returns the HTML-embedded description to any
client, unescaped script tags execute in the victim's browser.

\textbf{VULN-11 — Command Injection (CWE-77).}
The admin ping utility passes user-supplied hostnames directly to the
Node.js \texttt{child\_process.exec} function:
\begin{lstlisting}[language=JavaScript,caption=Vulnerable command execution]
exec(`ping -c 4 ${req.query.host}`, (err, stdout) =>
  res.json({ output: stdout }));
\end{lstlisting}
An attacker can execute arbitrary shell commands by appending a semicolon
and a payload (e.g., \texttt{localhost;cat /etc/passwd}).

\textbf{VULN-13 — Insecure Dockerfile.}
The container is built from an end-of-life base image (\texttt{node:14-alpine})
and the application runs as the container root user.
Trivy reports multiple CVEs against the outdated base image, and NIST
guidelines explicitly recommend running containers as non-root
principals~\cite{nistcontainer}.

\subsection{Dependency Vulnerabilities}

VulnAPI intentionally pins outdated dependency versions to produce a
non-trivial npm audit report.
Table~\ref{tab:deps} lists the pinned packages and their known CVEs.

\begin{table}[t]
  \caption{Pinned Vulnerable Dependencies}
  \label{tab:deps}
  \centering
  \scriptsize
  \begin{tabular}{@{}lll@{}}
    \toprule
    Package & Version & Notable CVEs \\
    \midrule
    express      & 4.17.1 & CVE-2022-24999 (ReDoS) \\
    jsonwebtoken & 8.5.1  & CVE-2022-23529 (secret forgery) \\
    lodash       & 4.17.4 & CVE-2019-10744 (prototype pollution) \\
    marked       & 0.3.6  & CVE-2022-21681 (ReDoS) \\
    node-fetch   & 2.6.0  & CVE-2022-0235 (SSRF redirect) \\
    body-parser  & 1.19.0 & CVE-2024-45590 (DoS) \\
    \bottomrule
  \end{tabular}
\end{table}

% ── V. DEVSECOPS PIPELINE ──────────────────────────────────────────
\section{DevSecOps CI/CD Pipeline}

\subsection{Pipeline Overview}

The pipeline is defined as a single GitHub Actions workflow file
(\texttt{.github/workflows/devsecops.yml}).
It is triggered on every push to any branch and on pull requests targeting
\texttt{main}.
Six jobs execute in a DAG (directed acyclic graph) where Jobs~3--5 depend on
Job~1's Docker image artifact, and Job~6 depends on all preceding jobs but
runs unconditionally (\texttt{if: always()}) to ensure the dashboard is always
refreshed even when scans detect failures.

\subsection{Job 1: Dependency Scan (npm audit)}

The dependency scan job checks out the repository, installs Node.js~18, and
runs \texttt{npm audit -{}-json}, capturing the full advisory tree in
\texttt{reports/npm-audit.json}.
The pipeline uses \texttt{|| true} to prevent a non-zero npm audit exit code
from blocking subsequent stages—the philosophy is to \emph{report}, not
\emph{block}, because blocking on every advisory in a legacy codebase is
counterproductive.
The JSON report is uploaded as a GitHub Actions artifact with a 30-day
retention window.

\subsection{Job 2: Static Application Security Testing (SonarCloud)}

The SAST job uses the official \texttt{SonarSource/sonarqube-scan-action}
with a full Git history checkout (\texttt{fetch-depth: 0}) to enable
SonarCloud's new-code and blame analysis features.
Configuration is provided via \texttt{sonar-project.properties}, specifying
the project key, organization, and source roots.
After the scan completes, the job queries the SonarCloud Issues REST API
(\texttt{/api/issues/search}) to export all open, unresolved findings
as a structured JSON file that the dashboard can consume.
This two-phase approach (scan + export) decouples issue ingestion from the
SonarCloud web UI, enabling the dashboard to surface findings without
requiring users to log into a separate tool.

\subsection{Job 3: Docker Image Build}

The Docker build job creates a tagged image
(\texttt{vulnerable-api:\$\{github.sha\}}) and serializes it to a gzipped
tarball via \texttt{docker save | gzip}.
The tarball is uploaded as a short-lived artifact (1-day retention) that is
shared with the Trivy and DAST jobs, ensuring both scan the exact same image
that would be deployed to production.
Using the commit SHA as the image tag guarantees immutability and full
traceability.

\subsection{Job 4: Container Scan (Trivy)}

The Trivy job downloads the image tarball, loads it, and runs two Trivy
passes: one producing SARIF output uploaded to the GitHub Security
Code Scanning tab, and one producing JSON output consumed by the dashboard.
Both passes cover all severity levels (CRITICAL, HIGH, MEDIUM, LOW) with
\texttt{exit-code: 0} to prevent build blocking.
Uploading SARIF to the Security tab provides an always-on, per-commit
vulnerability history directly inside the GitHub repository interface.

\subsection{Job 5: Dynamic Application Security Testing (OWASP ZAP)}

The DAST job presents a non-trivial networking challenge: OWASP ZAP runs
inside its own Docker container with \texttt{-{}-network host}, while the
Node.js application under test runs directly on the GitHub Actions runner.
The standard \texttt{localhost} address is inaccessible from within the ZAP
container because it refers to the container's own loopback interface.

The solution resolves the Docker bridge gateway IP at runtime:
\begin{lstlisting}[language=bash,caption=Runtime gateway resolution]
GATEWAY=$(ip route show default \
  | awk '/default/ {print $3}' | head -1)
echo "ZAP_TARGET=http://${GATEWAY:-172.17.0.1}:3000" \
  >> "$GITHUB_ENV"
\end{lstlisting}
ZAP then performs a full active scan against that IP, exercising all five
route groups.
The job includes a 30-attempt health check polling loop to ensure the
application is ready before ZAP begins, with a two-second interval between
attempts.

\subsection{Job 6: Dashboard Build and Deploy}

The final job downloads all JSON reports into \texttt{dashboard/public/reports/},
generates a \texttt{pipeline-meta.json} file containing the run ID, commit
SHA, branch, actor, timestamp, and per-job status strings, installs
dependencies, and builds the Vite/React project.
The build output is not published to GitHub Pages; instead, a secondary step
commits the updated JSON reports back to \texttt{main} tagged
\texttt{[skip ci]} to prevent an infinite pipeline loop.
Render detects this push and automatically redeploys the static dashboard,
ensuring the live site always reflects the most recent scan results.
PRs perform the build as a dry-run check without triggering a Render deploy.

% ── VI. SECURITY INTELLIGENCE DASHBOARD ───────────────────────────
\section{Security Intelligence Dashboard}

\subsection{Technology Stack}

The dashboard is built with React~18 and Vite, deployed as a static single-page
application (SPA) with no server-side component.
Recharts provides the interactive bar chart visualizations.
All data is fetched at runtime via the browser's Fetch API from relative paths
resolved against the Render deployment base URL, avoiding any hard-coded origin
dependencies.
A custom \texttt{useReport} hook encapsulates loading state, error handling,
and JSON parsing for each report file.

\subsection{Dashboard Views}

The application provides five views accessible via a persistent left sidebar:

\textbf{Command Center (Overview).}
The landing view aggregates findings from all four scanners into summary
statistic cards (total findings, critical, high, medium, low), a pipeline
stage status strip showing per-job pass/fail/skipped indicators, a grouped
bar chart of findings by tool and severity, and a prioritized Top Threats
panel listing the twelve highest-severity findings across all tools with
their source tool, CWE location, and severity badge.

\textbf{Dependency Scan.}
Renders the npm audit JSON in a searchable, sortable table showing each
vulnerable package, its installed version, fixed version, severity, and
advisory URL.

\textbf{Static Analysis.}
Renders SonarCloud issues grouped by rule and severity.
Each row links to the affected file and line number, enabling one-click
navigation to the offending code.

\textbf{Infrastructure.}
Renders Trivy results organized by OS package and application library,
with CVE identifiers, CVSS scores, installed versions, and fix versions
where available.

\textbf{Live Probe.}
Renders OWASP ZAP alerts sorted by risk code, including evidence snippets,
solution guidance, and the affected URI.

\subsection{Data Flow}

The dashboard is designed to be decoupled from the pipeline.
JSON report files are static assets served alongside the SPA.
On each page visit, the browser fetches the latest reports directly from
Render's CDN infrastructure, meaning the dashboard always reflects the
most recent pipeline run without any API server or database.
The \texttt{pipeline-meta.json} file provides the last-scanned timestamp and
branch/commit context rendered in the top bar.

% ── VII. REMEDIATION (app-fixed) ───────────────────────────────────
\section{Remediation Strategy}

The \texttt{app-fixed} directory contains a fully remediated version of VulnAPI
that demonstrates a clean pipeline run.
Table~\ref{tab:fixes} maps each vulnerability to its fix strategy.

\begin{table}[t]
  \caption{Vulnerability Remediation Summary}
  \label{tab:fixes}
  \centering
  \scriptsize
  \begin{tabular}{@{}lp{5.5cm}@{}}
    \toprule
    ID(s) & Remediation Applied \\
    \midrule
    VULN-01/02 & Secrets loaded from environment variables via \texttt{dotenv};
                 hard-coded values removed entirely. \\
    VULN-03    & CORS configured with an explicit origin allowlist from
                 \texttt{CORS\_ORIGINS} env var; wildcard origin removed. \\
    VULN-04/08 & All SQL queries converted to parameterized statements;
                 user input never concatenated into query strings. \\
    VULN-05    & Generic \texttt{"Internal server error"} message returned on
                 database errors; query text never included in responses. \\
    VULN-06    & Ownership check added: \texttt{req.user.id} compared against
                 the requested resource before returning data. \\
    VULN-07    & Password hash column excluded from all SELECT projections;
                 API responses return only non-sensitive fields. \\
    VULN-09    & Product descriptions HTML-escaped using \texttt{he} library
                 before storage and on output. \\
    VULN-10    & Path sanitized with \texttt{path.resolve} and validated to
                 remain within the allowed base directory. \\
    VULN-11    & Shell execution replaced with \texttt{child\_process.execFile}
                 using a validated hostname argument, eliminating shell
                 interpretation. \\
    VULN-12    & Admin routes protected by JWT verification middleware;
                 role checked for \texttt{admin} before allowing access. \\
    VULN-13    & Base image updated to \texttt{node:18-alpine}; application runs
                 as a dedicated non-root \texttt{node} user. \\
    \bottomrule
  \end{tabular}
\end{table}

\subsection{Additional Hardening}

Beyond the thirteen targeted fixes, \texttt{app-fixed} applies defense-in-depth
measures.
The \texttt{helmet} middleware adds nine security HTTP headers including
\texttt{Content-Security-Policy}, \texttt{X-Frame-Options}, and
\texttt{X-Content-Type-Options}.
Request body sizes are capped at 10 KB to mitigate denial-of-service through
oversized payloads.
A \texttt{express-rate-limit} middleware restricts the login endpoint to ten
attempts per fifteen-minute window, providing brute-force protection.
All dependency versions are updated to their latest non-breaking releases to
eliminate the CVEs listed in Table~\ref{tab:deps}.

% ── VIII. RESULTS AND EVALUATION ───────────────────────────────────
\section{Results and Evaluation}

\subsection{Vulnerability Detection Coverage}

Table~\ref{tab:coverage} shows which pipeline tools detected each of the
thirteen intentional vulnerabilities.

\begin{table}[t]
  \caption{Vulnerability Detection Coverage by Tool}
  \label{tab:coverage}
  \centering
  \scriptsize
  \begin{tabular}{@{}lcccc@{}}
    \toprule
    Vuln ID & npm audit & SonarCloud & Trivy & ZAP \\
    \midrule
    VULN-01 (JWT secret)       & --  & \checkmark & --  & --  \\
    VULN-02 (hard-coded creds) & --  & \checkmark & --  & --  \\
    VULN-03 (CORS)             & --  & \checkmark & --  & \checkmark \\
    VULN-04 (SQL inj. login)   & --  & \checkmark & --  & \checkmark \\
    VULN-05 (SQL error leak)   & --  & \checkmark & --  & \checkmark \\
    VULN-06 (IDOR)             & --  & --  & --  & \checkmark \\
    VULN-07 (password leak)    & --  & \checkmark & --  & \checkmark \\
    VULN-08 (SQL inj. search)  & --  & \checkmark & --  & \checkmark \\
    VULN-09 (stored XSS)       & --  & \checkmark & --  & \checkmark \\
    VULN-10 (path traversal)   & --  & \checkmark & --  & \checkmark \\
    VULN-11 (cmd injection)    & --  & \checkmark & --  & \checkmark \\
    VULN-12 (no auth)          & --  & --  & --  & \checkmark \\
    VULN-13 (root Dockerfile)  & --  & --  & \checkmark & --  \\
    Dep.\ CVEs (6 packages)   & \checkmark & --  & \checkmark & --  \\
    \bottomrule
  \end{tabular}
\end{table}

All thirteen intentional vulnerabilities are detected by at least one pipeline
tool.
SonarCloud achieves the broadest code-level coverage, detecting ten of
thirteen vulnerabilities through static analysis.
OWASP ZAP complements it by dynamically confirming exploitability at runtime
for nine vulnerabilities.
Trivy provides unique value for infrastructure-level issues such as VULN-13
and OS-level CVEs in the base image.
npm audit catches all six dependency CVEs, including transitive ones not
visible to static analysis.

The tools are complementary rather than redundant.
No single tool detects all vulnerabilities, validating the multi-layer
architecture.

\subsection{Pipeline Execution Metrics}

Table~\ref{tab:timing} presents typical execution times observed across
multiple pipeline runs.

\begin{table}[t]
  \caption{Typical Job Execution Times}
  \label{tab:timing}
  \centering
  \scriptsize
  \begin{tabular}{@{}lrl@{}}
    \toprule
    Job & Duration & Notes \\
    \midrule
    Dependency Scan (npm audit) &  \textasciitilde{}45 s  & Includes \texttt{npm ci} \\
    SAST (SonarCloud)           &  \textasciitilde{}2 min & Full Git history fetch \\
    Docker Build                &  \textasciitilde{}90 s  & Layer caching not used \\
    Container Scan (Trivy)      &  \textasciitilde{}3 min & DB download + scan \\
    DAST (OWASP ZAP)            &  \textasciitilde{}8 min & Full active scan \\
    Dashboard Build \& Deploy   &  \textasciitilde{}2 min & Includes npm build \\
    \midrule
    \textbf{Total (wall clock)} & \textasciitilde{}12 min & Jobs 2--5 run in parallel \\
    \bottomrule
  \end{tabular}
\end{table}

Because Jobs~2 through~5 execute in parallel, the effective wall-clock time
is dominated by the slowest job (ZAP) rather than the sum.
Total end-to-end pipeline time averages approximately twelve minutes on
GitHub Actions' standard Ubuntu runners.

\subsection{Fixed Application Results}

When the pipeline is run against \texttt{app-fixed}, all security findings
are resolved:
npm audit reports zero vulnerabilities, SonarCloud reports zero open issues,
Trivy reports zero CRITICAL/HIGH CVEs against the updated base image, and
OWASP ZAP reports only low-severity informational alerts.
The dashboard's Command Center shows a green ``Deploy Unblocked'' indicator,
validating the end-to-end workflow.

\subsection{Dashboard Performance}

The React SPA loads the landing page in under two seconds on a broadband
connection, including all four JSON fetch calls.
Report files range from approximately 15 KB (SonarCloud with few issues) to
over 2 MB (Trivy with a full CVE database match), so the dashboard renders
scan-specific components lazily after their data arrives.

% ── IX. DISCUSSION ─────────────────────────────────────────────────
\section{Discussion}

\subsection{Tool Complementarity}

A key finding is that no single scanning tool is sufficient.
SonarCloud excels at detecting code-level issues such as SQL injection and
hard-coded secrets because it can trace data flows and recognize anti-patterns
at the source level.
However, it cannot distinguish between a vulnerability that is theoretically
present and one that is actually exploitable in a running deployment.
ZAP fills this gap through dynamic probing.
Trivy adds a dimension neither SAST nor DAST addresses: the security posture
of the runtime environment itself—the OS, package manager, and base image.
npm audit addresses the supply chain attack surface, which is absent from all
other tools.
This complementary coverage justifies maintaining all four tools in the
pipeline despite increased complexity and runtime.

\subsection{Developer Experience}

The system is designed to minimize friction.
The only setup requirement is adding a single \texttt{SONAR\_TOKEN} secret to
the GitHub repository.
All other tools are free and require no account registration.
Results are surfaced in three places developers already monitor: the
GitHub Actions UI (per-job logs), the GitHub Security tab (SARIF), and the
deployed dashboard.
This redundancy is intentional: different team members consult different
surfaces, and consistent data across all of them reduces the chance that a
critical finding is missed.

\subsection{Limitations}

\textbf{SQLite in-process database.}
VulnAPI uses SQLite running in the same process as the Express server.
In production deployments, the database would be a separate service with
network isolation.
ZAP's network-based attacks would not change substantially, but the Trivy
scan would need to cover the database container image as well.

\textbf{Authenticated DAST.}
The current ZAP configuration uses an unauthenticated full scan.
Endpoints behind JWT authentication (e.g., the \texttt{/api/users} route) may
not be fully exercised.
Future work should configure ZAP with a valid JWT to enable authenticated
active scanning.

\textbf{False negative risk in SAST.}
SonarCloud's JavaScript analysis does not yet support all taint-tracking
scenarios.
Complex multi-hop data flows where user input passes through multiple function
calls before reaching a sink may be missed.

\textbf{Pipeline execution time.}
At twelve minutes per push, the pipeline adds meaningful latency to the
developer feedback loop.
Caching Trivy's vulnerability database and enabling Docker layer caching
could reduce this to approximately six minutes without sacrificing coverage.

\subsection{Future Work}

Several extensions are planned.
First, integrating GitHub Dependabot to automatically open pull requests for
outdated dependencies would complement the passive reporting of npm audit.
Second, adding a Software Bill of Materials (SBOM) generation step using
Syft would provide a machine-readable inventory of all components for
compliance purposes.
Third, extending the dashboard to track finding trends across commits would
enable teams to monitor their security debt trajectory over time.
Finally, adding infrastructure-as-code scanning (e.g., Checkov or tfsec) to
the pipeline would extend security coverage to Terraform and Kubernetes
manifests.

% ── X. CONCLUSION ──────────────────────────────────────────────────
\section{Conclusion}

This paper presented SecureShift, a fully automated DevSecOps pipeline that
integrates dependency scanning, SAST, container scanning, and DAST into a
single GitHub Actions workflow.
The pipeline was evaluated against VulnAPI, an intentionally vulnerable
Node.js/Express application containing thirteen documented weaknesses across
seven OWASP Top~10 categories.
The four-tool combination detected all thirteen vulnerabilities with no manual
configuration beyond a single API token.
The accompanying React Security Intelligence Dashboard aggregates scan output
into a unified, real-time view accessible to all stakeholders without
specialized tooling.

The remediated \texttt{app-fixed} application confirmed that the pipeline
correctly transitions from a failing to a passing state once all vulnerabilities
are addressed, validating the end-to-end workflow.

SecureShift demonstrates that comprehensive security scanning is achievable at
zero marginal infrastructure cost using freely available tools and GitHub's
native CI/CD platform.
The complete implementation is available as an open-source reference at
\texttt{github.com/neelp03/CMPE272-project} to serve as a practical starting
point for teams adopting DevSecOps practices.

% ── REFERENCES ─────────────────────────────────────────────────────
\begin{thebibliography}{00}

\bibitem{nist2002cost}
R.~Pressman, \emph{Software Engineering: A Practitioner's Approach},
7th~ed.\ New York, NY, USA: McGraw-Hill, 2010.

\bibitem{devsecops2019}
S.~Hasan, Z.~Ali, S.~Madnick, and S.~Velez-Caicedo,
``Understanding How DevSecOps Addresses the Security Gap in DevOps,''
in \emph{Proc.\ IEEE Int.\ Conf.\ Softw.\ Archit.\ Companion (ICSA-C)},
Hamburg, Germany, 2019, pp.\ 252--259.

\bibitem{myrbakken2017}
H.~Myrbakken and R.~Colomo-Palacios,
``DevSecOps: A Multivocal Literature Review,''
in \emph{Proc.\ Int.\ Conf.\ Software Process Improvement and Capability
Determination (SPICE)}, Palma, Spain, 2017, pp.\ 17--29.

\bibitem{rajapakse2022}
V.~N.~Rajapakse, M.~Zahedi, M.~A.~Babar, and H.~Shen,
``Challenges and Solutions When Adopting DevSecOps: A Systematic Review,''
\emph{Inf.\ Softw.\ Technol.}, vol.~141, p.\ 106700, 2022.

\bibitem{npmaudit}
npm, Inc., ``npm audit,'' npm Documentation, 2023. [Online].
Available: \url{https://docs.npmjs.com/cli/commands/npm-audit}

\bibitem{sonarcloud}
SonarSource, ``SonarCloud Documentation,'' 2024. [Online].
Available: \url{https://docs.sonarcloud.io}

\bibitem{trivy2021}
A.~Security, ``Trivy: A Comprehensive and Versatile Security Scanner,''
Aqua Security, GitHub, 2021. [Online].
Available: \url{https://github.com/aquasecurity/trivy}

\bibitem{zapproxy}
OWASP Foundation, ``OWASP Zed Attack Proxy (ZAP),'' OWASP, 2024. [Online].
Available: \url{https://www.zaproxy.org}

\bibitem{siem2020}
K.~Kent and M.~Souppaya, ``Guide to Computer Security Log Management,''
NIST Special Publication 800-92, National Institute of Standards and Technology,
Gaithersburg, MD, USA, 2006.

\bibitem{perez2021}
R.~Pérez-Castillo, I.~G.~Rodríguez de Guzmán, and M.~Piattini,
``Security Smells in Ansible and Chef Scripts: A Replication Study,''
\emph{ACM Trans.\ Softw.\ Eng.\ Methodol.}, vol.\ 30, no.\ 1, pp.\ 1--34,
2021.

\bibitem{nistcontainer}
M.~Souppaya and K.~Scarfone, ``Application Container Security Guide,''
NIST Special Publication 800-190, National Institute of Standards and Technology,
Gaithersburg, MD, USA, 2017.

\end{thebibliography}

\end{document}
